% ---
% title: A game with an urn
% tags: random-variables, distribution-function, expectation, variance
% ---

\begin{problem}
	Two balls are drawn at random, without replacement, from an urn that
	contains 8 white, 4 black and 2 orange balls. We win 2 euros for each black
	ball drawn and lose 1 euro for each white ball drawn; an orange ball brings
	nothing. Let $X$ be our net winnings (in euros).
	(a) Find the probability mass function of $X$.
	(b) Find the distribution function $F(x)=P(X\le x)$ and sketch its graph.
	(c) Find $P(X>0)$ and $F(1)$.
	(d) Find $\E[X]$, $\Var(X)$ and $\SD(X)$.
\end{problem}

\begin{slideproblem}
	Two balls are drawn without replacement from an urn with 8 white, 4 black
	and 2 orange balls. Each black ball wins 2 euros, each white ball loses
	1 euro, orange: nothing. $X$ = net winnings.
	(a) pmf of $X$; \ (b) distribution function $F(x)=P(X\le x)$, with a sketch;
	\ (c) $P(X>0)$ and $F(1)$; \ (d) $\E[X]$, $\Var(X)$, $\SD(X)$.
\end{slideproblem}

\begin{slides}
	\small
	\exstatement
	\pause
	\sol (a) All pairs of balls are equally likely:
	$\binom{14}{2}=\frac{14\cdot13}{2}=91$ pairs.
	\pause
	List the possible pairs by colors, count them, and find the winnings:
	\begin{center}
		\renewcommand{\arraystretch}{1.15}
		\begin{tabular}{l|cccccc|c}
			colors & WW & WO & OO & WB & BO & BB & total \\
			\hline
			number of pairs & $\binom82=28$ & $8\cdot2=16$ & $1$ & $8\cdot4=32$
			& $4\cdot2=8$ & $\binom42=6$ & $91$ \\
			\pause
			value $x$ of $X$ & $-2$ & $-1$ & $0$ & $1$ & $2$ & $4$ & \\
			\pause
			$p(x)=P(X=x)$ & $\frac{28}{91}$ & $\frac{16}{91}$ & $\frac{1}{91}$
			& $\frac{32}{91}$ & $\frac{8}{91}$ & $\frac{6}{91}$ & $1$ \\
		\end{tabular}
	\end{center}
	\pause
	Check: the probabilities add up to $\frac{91}{91}=1$.
\end{slides}

\begin{slides}
	\small
	\begin{columns}[T]
		\begin{column}{.42\textwidth}
			(b) Add up the pmf from the left:
			\[ F(x)=\begin{cases}
					0,                & x<-2,         \\[.5mm]
					\frac{28}{91},    & -2\le x<-1,   \\[.5mm]
					\frac{44}{91},    & -1\le x<0,    \\[.5mm]
					\frac{45}{91},    & 0\le x<1,     \\[.5mm]
					\frac{77}{91},    & 1\le x<2,     \\[.5mm]
					\frac{85}{91},    & 2\le x<4,     \\[.5mm]
					1,                & x\ge 4.
				\end{cases} \]
		\end{column}
		\begin{column}{.56\textwidth}
			\centering
			\begin{tikzpicture}[xscale=.62,yscale=3.4]
				\draw[->] (-3.4,0) -- (5.4,0) node[right] {\scriptsize $x$};
				\draw[->] (0,-.05) -- (0,1.15) node[above] {\scriptsize $F(x)$};
				\foreach \x in {-3,-2,-1,1,2,3,4,5}
				\draw (\x,.02) -- (\x,-.02) node[below] {\tiny $\x$};
				\draw (.06,1) -- (-.06,1) node[left] {\tiny $1$};
				\draw[navy,thick] (-3.4,0) -- (-2,0);
				\foreach \a/\b/\v/\t in {-2/-1/28/{28},-1/0/44/{44},0/1/45/{45},
						1/2/77/{77},2/4/85/{85}}
					{
						\draw[navy,thick] (\a,\v/91) -- (\b,\v/91);
						\fill[navy] (\a,\v/91) circle (.07 and .013);
						\draw[navy,fill=white] (\b,\v/91) circle (.07 and .013);
					}
				\draw[navy,thick] (4,1) -- (5.4,1);
				\fill[navy] (4,1) circle (.07 and .013);
				\node[above,teal] at (-1.5,28/91) {\tiny $\frac{28}{91}$};
				\node[above,teal] at (-.5,44/91+.04) {\tiny $\frac{44}{91}$};
				\node[above,teal] at (.55,45/91) {\tiny $\frac{45}{91}$};
				\node[above,teal] at (1.5,77/91) {\tiny $\frac{77}{91}$};
				\node[above,teal] at (3,85/91) {\tiny $\frac{85}{91}$};
			\end{tikzpicture}
		\end{column}
	\end{columns}
	\smallskip
	Right-continuous step function; the jump at $x$ is $p(x)$; filled dot = the
	value at the jump point.
	\pause
	\smallskip

	(c) $P(X>0)=1-F(0)=1-\frac{45}{91}=\frac{46}{91}\approx0.505$
	\ (or $\frac{32+8+6}{91}$).\\
	\phantom{(c) }$F(1)=P(X\le1)=\frac{77}{91}=\frac{11}{13}\approx0.846$.
\end{slides}

\begin{slides}
	\small
	(d) Expected value: each value times its probability, then add.
	\[ \E[X]=\frac{(-2)\cdot28+(-1)\cdot16+0\cdot1+1\cdot32+2\cdot8+4\cdot6}{91}
		=\frac{-56-16+0+32+16+24}{91}=0. \]
	On average we neither win nor lose: a \alert{fair game}.
	\pause
	\medskip

	Variance: $\Var(X)=\E[X^2]-(\E[X])^2$. First $\E[X^2]$: square the values,
	keep the probabilities.
	\[ \E[X^2]=\frac{4\cdot28+1\cdot16+0\cdot1+1\cdot32+4\cdot8+16\cdot6}{91}
		=\frac{112+16+0+32+32+96}{91}=\frac{288}{91}. \]
	\pause
	\[ \Var(X)=\frac{288}{91}-0^2=\frac{288}{91}\approx3.165, \qquad
		\SD(X)=\sqrt{\frac{288}{91}}\approx1.779 \text{ euros}. \]
	\pause
	The game is fair, but risky: a typical deviation from 0 is about
	1.78 euros.
\end{slides}

\begin{solution}
	\idea Every pair of balls is equally likely. Group the pairs by their colors,
	count each group, and read off the winnings. The distribution function adds
	up the probabilities from the left; $\E[X]$ and $\E[X^2]$ are weighted
	averages with the probabilities as weights.

	\step{(a)} The urn has $8+4+2=14$ balls. The number of ways to choose 2 of
	them (order does not matter) is
	\[ \binom{14}{2}=\frac{14\cdot13}{2}=91, \]
	and each of these 91 pairs is equally likely. Now list the possible color
	combinations. For each one, count the pairs and compute the winnings
	($+2$ for black, $-1$ for white, $0$ for orange):
	\begin{itemize}
		\item two white (WW): $\binom82=28$ pairs, $X=-1-1=-2$;
		\item white and orange (WO): $8\cdot2=16$ pairs (one of 8 white balls
		      and one of 2 orange balls), $X=-1+0=-1$;
		\item two orange (OO): $\binom22=1$ pair, $X=0$;
		\item white and black (WB): $8\cdot4=32$ pairs, $X=-1+2=1$;
		\item black and orange (BO): $4\cdot2=8$ pairs, $X=2+0=2$;
		\item two black (BB): $\binom42=6$ pairs, $X=2+2=4$.
	\end{itemize}
	Check: $28+16+1+32+8+6=91$, so no pair is missing. Different color
	combinations give different values, so each probability is the number of
	pairs divided by 91:
	\begin{center}
		\renewcommand{\arraystretch}{1.3}
		\begin{tabular}{c|cccccc}
			$x$      & $-2$ & $-1$ & $0$ & $1$ & $2$ & $4$ \\
			\hline
			$p(x)$   & $\frac{28}{91}$ & $\frac{16}{91}$ & $\frac{1}{91}$
			         & $\frac{32}{91}$ & $\frac{8}{91}$ & $\frac{6}{91}$
		\end{tabular}
	\end{center}

	\step{(b)} $F(x)=P(X\le x)$ is the sum of $p(k)$ over all values
	$k\le x$. Below $-2$ there is no value, so $F(x)=0$. At each value of $X$
	the function jumps up by the probability of that value, and between two
	values it stays constant. Adding the probabilities from the left:
	$28$, $28+16=44$, $44+1=45$, $45+32=77$, $77+8=85$, $85+6=91$ (over 91).
	So
	\[ F(x)=\begin{cases}
			0,             & x<-2,       \\
			28/91,         & -2\le x<-1, \\
			44/91,         & -1\le x<0,  \\
			45/91,         & 0\le x<1,   \\
			77/91,         & 1\le x<2,   \\
			85/91,         & 2\le x<4,   \\
			1,             & x\ge 4.
		\end{cases} \]
	The graph is a step function. Each step includes its left endpoint (filled
	dot) but not its right endpoint (open circle): for example $F(-2)=28/91$,
	because $X\le-2$ includes $X=-2$. This is why a distribution function is
	right-continuous.
	\begin{center}
		\begin{tikzpicture}[xscale=.8,yscale=3.4]
			\draw[->] (-3.4,0) -- (5.4,0) node[right] {\small $x$};
			\draw[->] (0,-.05) -- (0,1.15) node[above] {\small $F(x)$};
			\foreach \x in {-3,-2,-1,1,2,3,4,5}
			\draw (\x,.02) -- (\x,-.02) node[below] {\scriptsize $\x$};
			\draw (.06,1) -- (-.06,1) node[left] {\scriptsize $1$};
			\draw[navy,thick] (-3.4,0) -- (-2,0);
			\foreach \a/\b/\v in {-2/-1/28,-1/0/44,0/1/45,1/2/77,2/4/85}
				{
					\draw[navy,thick] (\a,\v/91) -- (\b,\v/91);
					\fill[navy] (\a,\v/91) circle (.06 and .014);
					\draw[navy,fill=white] (\b,\v/91) circle (.06 and .014);
				}
			\draw[navy,thick] (4,1) -- (5.4,1);
			\fill[navy] (4,1) circle (.06 and .014);
			\node[above,teal] at (-1.5,28/91) {\scriptsize $\frac{28}{91}$};
			\node[above,teal] at (-.5,44/91+.04) {\scriptsize $\frac{44}{91}$};
			\node[above,teal] at (.55,45/91) {\scriptsize $\frac{45}{91}$};
			\node[above,teal] at (1.5,77/91) {\scriptsize $\frac{77}{91}$};
			\node[above,teal] at (3,85/91) {\scriptsize $\frac{85}{91}$};
		\end{tikzpicture}
	\end{center}

	\step{(c)} ``$X>0$'' is the complement of ``$X\le0$'', so
	\[ P(X>0)=1-F(0)=1-\frac{45}{91}=\frac{46}{91}\approx0.505. \]
	Directly: the positive values are 1, 2 and 4, and
	$p(1)+p(2)+p(4)=\frac{32+8+6}{91}=\frac{46}{91}$. Next, $F(1)$ is read off
	the step that starts at 1:
	\[ F(1)=P(X\le1)=\frac{77}{91}=\frac{11}{13}\approx0.846. \]
	Note that $F(1)$ includes $X=1$; the probability $P(X<1)$ would be only
	$45/91$.

	\step{(d)} The expected value is the sum of the values, each multiplied by
	its probability:
	\[ \begin{aligned}
			\E[X] & =\frac{(-2)\cdot28+(-1)\cdot16+0\cdot1+1\cdot32+2\cdot8+4\cdot6}{91} \\
			      & =\frac{-56-16+0+32+16+24}{91}=\frac{0}{91}=0.
		\end{aligned} \]
	On average we neither win nor lose: the game is fair.

	For the variance we use $\Var(X)=\E[X^2]-(\E[X])^2$. To get $\E[X^2]$,
	square each value and keep the same probabilities:
	\[ \begin{aligned}
			\E[X^2] & =\frac{(-2)^2\cdot28+(-1)^2\cdot16+0^2\cdot1+1^2\cdot32
			+2^2\cdot8+4^2\cdot6}{91} \\
			        & =\frac{112+16+0+32+32+96}{91}=\frac{288}{91}.
		\end{aligned} \]
	Hence
	\[ \Var(X)=\frac{288}{91}-0^2=\frac{288}{91}\approx3.165,\qquad
		\SD(X)=\sqrt{\frac{288}{91}}\approx1.779 \text{ euros}. \]
	The variance is measured in squared euros; the standard deviation is back
	in euros, so it is the number to interpret: a typical deviation of the
	winnings from their mean 0 is about 1.78 euros.
	\begin{answer}
		(a) $p(-2)=\frac{28}{91}$, $p(-1)=\frac{16}{91}$, $p(0)=\frac{1}{91}$,
		$p(1)=\frac{32}{91}$, $p(2)=\frac{8}{91}$, $p(4)=\frac{6}{91}$
		\quad (b) step function with values $0$, $\frac{28}{91}$, $\frac{44}{91}$,
		$\frac{45}{91}$, $\frac{77}{91}$, $\frac{85}{91}$, $1$ (jumps at
		$-2,-1,0,1,2,4$)
		\quad (c) $P(X>0)=\frac{46}{91}\approx0.505$, $F(1)=\frac{11}{13}\approx0.846$
		\quad (d) $\E[X]=0$, $\Var(X)=\frac{288}{91}\approx3.165$,
		$\SD(X)\approx1.779$
	\end{answer}
	\begin{remark}
		\textbf{Common mistake:} computing $\Var(X)$ as $\E[X^2]$ without
		subtracting $(\E[X])^2$. Here it makes no difference only because
		$\E[X]=0$. Another common mistake in (b) is to put the filled dot at the
		right end of a step: $F$ jumps \emph{at} the value, so $F(x)$ already
		contains $p(x)$.
	\end{remark}
\end{solution}
